%%%%%%%%%%%%%%%%%%%%%%%%%%%%%%%%%%%%%%%%%%%%%%%%%%%%%%%%%%%%%%%%%%%%%%
% How to use writeLaTeX: 
%
% You edit the source code here on the left, and the preview on the
% right shows you the result within a few seconds.
%
% Bookmark this page and share the URL with your co-authors. They can
% edit at the same time!
%
% You can upload figures, bibliographies, custom classes and
% styles using the files menu.
%
%%%%%%%%%%%%%%%%%%%%%%%%%%%%%%%%%%%%%%%%%%%%%%%%%%%%%%%%%%%%%%%%%%%%%%

\documentclass[12pt]{article}

\usepackage{sbc-template}

\usepackage{graphicx,url}

%\usepackage[brazil]{babel}   
\usepackage[utf8]{inputenc}  

     
\sloppy

\title{Exercício 5 Sistemas Distribuídos}

\author{Christian Aguiar Plentz\inst{1}, João Roberto Lemos Fidellis\inst{2} }


\address{Ciência da Computação -- Universidade do Vale do Rio dos Sinos (UNISINOS)\\
  Caixa Postal 275 -- 93022-750 -- São Leopoldo -- RS -- Brazil
  \email{\{plentzchristian,JRFIDELLIS\}@edu.unisinos.br}
}


\begin{document} 

\maketitle

 \section{Escreva o passo a passo Alice-Bob para implementarmos o seguinte statement}
	\textit{Bob gostaria de enviar um documento assinado e com criptografia assimetrica para alice}

\subsection{Bob}

\texttt{
\noindent Recebe PubA de Alice \\
Cria as chaves: PubB e PrivB \\
Passa PubB para Alice \\
HASH = hash(sha1, documento()) \\
SIGN = sign(PrivB, HASH, sha1) \\
MSG = append(documento(), SIGN) \\
FINAL = cript(rsa, PubA, MSG) \\
send(alice, FINAL) \\
}

\subsection{Alice}
\texttt{
\noindent Cria as chaves: PubA e PrivA \\
Passa PubA para Bob \\
Recebe PubB de Bob \\
Recebe FINAL \\
READ = decrypt(rsa, PrivA, FINAL) \\
if (READ == NULL) \\
\indent PROBLEMA \\
else \\
\indent SIGN = pega\_final(READ) \\
\indent BODY = pega\_corpo(READ) \\
\indent HASH = hash(sha1, BODY) \\
\indent RESULT = verify(PubB, HASH, SIGN, sha1) \\
\indent if RESULT \\
\indent \indent PASS \\
\indent else \\
\indent	\indent PROBLEMA \\
}

\end{document}

